\section{Finite algebraic topology, sheaves, and Boolean algebras}
\label{sec:algebraic}

This section treats the index entries \emph{Betti Number}, \emph{Euler
Characteristic}, \emph{Genus}, \emph{Coloring}, the sheaf-related entries, and
\emph{Sigma-Algebra}.  The objects are finite tables or finite presentations.
Consequently the results below are uniform arithmetic theorems about those
presentations, not assertions that an arbitrary topological space has a
finite presentation.  The distinction matters especially for homotopy
invariance, surface classification, and countable unions.

\subsection{Chain complexes and exact computation of homology}

A finite rational chain complex consists of a natural number $d$, dimensions
$c_0,\ldots,c_d$, and rational matrices
\[
 D_k:\Q^{c_k}\longrightarrow\Q^{c_{k-1}}
 \qquad(1\leq k\leq d)
\]
whose consecutive products vanish.  Put $c_{-1}=c_{d+1}=0$ and let $D_0$ and
$D_{d+1}$ be the zero maps to and from the zero-dimensional spaces.  A matrix
is a rectangular list; a rational entry is an integer numerator and a
positive integer denominator.  All equalities in this subsection are
checked by cross multiplication.  In particular, no equality test for real
numbers occurs.  Write
\[
 Z_k=\ker D_k,\qquad B_k=\operatorname{im}D_{k+1},\qquad
 H_k=Z_k/B_k.
\]
The quotient notation will be implemented by a chosen basis and a projection,
as described below, rather than by a primitive quotient sort.

\begin{proposition}[A01: the simplicial boundary squares to zero]
\label{prop:boundary-square}
Let $K$ be a finite list of nonempty increasing vertex lists, closed under
nonempty sublists.  For $k>0$ define
\[
 \partial[v_0,\ldots,v_k]
   =\sum_{i=0}^{k}(-1)^i[v_0,\ldots,\widehat v_i,\ldots,v_k],
 \qquad \partial[v]=0.
\]
The resulting integer matrices satisfy $\partial^2=0$.
\end{proposition}
\begin{proof}
Fix two distinct deleted vertices with original indices $i<j$.  Deleting
$j$ first and then $i$ contributes sign $(-1)^{i+j}$ to the remaining
face.  Deleting $i$ first and then the vertex now in position $j-1$
contributes $(-1)^{i+j-1}$.  These contributions cancel.  Every term in the
double sum belongs to exactly this pair, indexed by $i<j$.  For $k=1$ the
second boundary is zero by definition.  This is a finite reindexing of
integer sums, formalizable by bounded induction on the two indices.
\end{proof}

\begin{proposition}[A02: homology with explicit normal forms]
\label{prop:effective-homology}
There is a primitive recursive algorithm computing integers $r_k$ and
$b_k$, rational matrices giving cycle and boundary bases, and mutually inverse
coordinate maps
\[
 H_k\cong\Q^{b_k},\qquad
 r_k=\operatorname{rank}D_k,\qquad
 b_k=c_k-r_k-r_{k+1}.
\]
In particular $r_k+r_{k+1}\leq c_k$.
\end{proposition}
\begin{proof}
Process the columns of a matrix from left to right.  At each step scan the
unused rows in their fixed order; if a nonzero entry exists, choose the first
one, exchange rows, divide the pivot row by that entry, and subtract its
appropriate multiples from every other row.  If there is no such entry,
move to the next column.  The procedure performs at most the number of
columns many stages, each with a bounded row scan.  Every row operation has
an explicitly recorded inverse.  The final reduced matrix therefore has
the same kernel as the original matrix and has $r$ pivot columns.  Its
nonpivot variables are free; assigning each in turn the value $1$ and the
others $0$, and solving the pivot equations, gives $c-r$ independent kernel
vectors spanning the kernel.  The original pivot columns form an image
basis: the invertible row transformation preserves their independence and
their spanning relations.  This proves rank--nullity without selecting an
unspecified basis.

Apply this construction to $D_k$ and $D_{k+1}$.  The chain identity says that
each of the $r_{k+1}$ image-basis vectors is a cycle.  Express these vectors
in the computed basis of $Z_k$, by the same elimination.  They remain
independent.  Starting with them, scan the cycle basis and append a vector
exactly when the enlarged list remains independent.  The resulting list is
a basis of $Z_k$ beginning with a basis of $B_k$.  Its remaining
$c_k-r_k-r_{k+1}$ vectors span a complement $L_k$, and
$Z_k=B_k\oplus L_k$.  Projection onto the last coordinates is the desired
normal form for a homology class.  Its kernel is precisely $B_k$, so it
also proves independence of the chosen cycle representative.

All scans have explicit finite bounds.  For a size bound, encode an
unreduced rational by a pair with absolute numerator and denominator at
most $S$.  Any addition, subtraction, multiplication, or division by a
tested nonzero rational produces a pair bounded by $2(S+1)^2$.
Starting from the largest input entry, iterate this bound for the number of
arithmetic operations of the fixed elimination schedule.  For a matrix
with at most $N$ rows and columns, $20(N+1)^4$ elementary arithmetic and
comparison operations is a sufficient loose bound for elimination,
recording row operations, and extracting bases.  A finite composition of
these schedules gives a primitive recursive bound for the entire complex.
Thus the proof is about bounded lists and identities of rational pairs.
\end{proof}

\begin{theorem}[A03: finite Euler--Poincar\'e formula]
\label{thm:euler-poincare}
For every finite rational chain complex,
\[
 \chi(C):=\sum_{k=0}^{d}(-1)^kc_k
       =\sum_{k=0}^{d}(-1)^kb_k.
\]
For a finite simplicial complex, $c_k$ is its number of $k$-simplices.
\end{theorem}
\begin{proof}
The preceding proposition gives $c_k=b_k+r_k+r_{k+1}$.  Substitute this in
the alternating sum.  The occurrence of $r_j$ from index $j$ has sign
$(-1)^j$; its occurrence from index $j-1$ has sign $(-1)^{j-1}$.  They
cancel for $1\leq j\leq d$, and the two endpoint terms are zero because
$r_0=r_{d+1}=0$.  This cancellation is a finite sum identity.  There is no
limit and no infinite-dimensional quotient.
\end{proof}

\begin{theorem}[A04: certified chain-homotopy invariance]
\label{thm:chain-homotopy}
A list of matrices $F_k:C_k\to C'_k$ satisfying
$D'_kF_k=F_{k-1}D_k$ induces computable linear maps on the normal forms
$H_k(C)\to H_k(C')$.  If
\[
 F_k-G_k=D'_{k+1}S_k+S_{k-1}D_k,
\]
then $F$ and $G$ induce the same homology map.  Consequently, supplied
chain maps $F,G$ and supplied chain homotopies $GF\simeq 1$ and
$FG\simeq 1$ imply equality of all Betti numbers and Euler characteristics.
\end{theorem}
\begin{proof}
For a cycle $z$, the commuting-square identity gives $D'Fz=FDz=0$.
For a boundary $D w$, it gives $F D w=D'F w$, a boundary.  Apply $F$ to
the chosen complement basis and project onto the target complement to
obtain its homology matrix.  Changes of representatives alter the answer
by a boundary and hence by zero after projection.  If $Dz=0$, the
homotopy equation reduces to $(F-G)z=D'Sz$, proving equality of the
induced maps.  Applying this to the two supplied homotopies yields inverse
homology matrices.  Independent columns and their inverse images then
give equal dimensions by finite rank--nullity.  The Euler conclusion
follows from Theorem~\ref{thm:euler-poincare}.
\end{proof}

This is an invariance theorem for \emph{certified chain homotopies}.
Turning an arbitrary continuous homotopy into such a certificate is a
separate theorem and is not smuggled into its hypothesis.

\begin{proposition}[A05: elementary collapse]
\label{prop:collapse}
Suppose $\sigma$ is a codimension-one face of $\tau$ in a finite
simplicial complex, and $\tau$ is the only simplex properly containing
$\sigma$.  Deleting $\sigma$ and $\tau$ gives a subcomplex $K'$ with the
same rational Betti numbers and Euler characteristic.
\end{proposition}
\begin{proof}
Write $\partial\tau=\varepsilon\sigma+b$, where
$\varepsilon\in\{1,-1\}$ and $b\in C(K')$.  Let $i:C(K')\to C(K)$ be
inclusion.  Define $r(\tau)=0$, $r(\sigma)=-\varepsilon b$, and let $r$
fix every other basis simplex.  Define $h(\sigma)=\varepsilon\tau$ and
let $h$ vanish on all other basis simplices.  Directly on $\sigma$ and
$\tau$ one checks
\[
 ir=1-\partial h-h\partial,\qquad ri=1.
\]
For any other simplex the identity holds because its boundary cannot
contain $\sigma$: that would make it another proper coface of $\sigma$.
Commuting $\partial$ with the displayed identity and using
$\partial^2=0$ shows that $r$ is a chain map; inclusion is visibly a
chain map.  These are the certificates required by
Theorem~\ref{thm:chain-homotopy}.  Alternatively the Euler characteristic
changes by deleting exactly one $k$-simplex and one $(k+1)$-simplex, whose
signed contributions sum to zero.  Repeating this argument for a supplied
finite collapse sequence preserves both conclusions.
\end{proof}

\begin{proposition}[A06: Euler additivity]
For finite subcomplexes $A,B$ of a common complex,
\[
 \chi(A\cup B)=\chi(A)+\chi(B)-\chi(A\cap B).
\]
Disjoint union adds each Betti number.
\end{proposition}
\begin{proof}
At each dimension, partition the simplex lists into those in $A\setminus
B$, those in $B\setminus A$, and those in $A\cap B$.  Counting this
partition proves the first assertion after taking alternating sums.  For
disjoint union, the boundary matrices are block diagonal.  Elimination
within each block gives the disjoint union of its cycle, boundary, and
complement bases, so the homology dimensions add.
\end{proof}

\begin{proposition}[A19: simplicial maps and the finite prism identity]
\label{prop:simplicial-prism}
An ordered vertex tuple denotes its oriented simplex: exchanging two
vertices changes its sign, and a tuple with repeated vertices denotes
zero.  A simplicial vertex map $f$ induces the chain map
$f_\#[v_0,\ldots,v_k]=[f(v_0),\ldots,f(v_k)]$.
If $f,g:K\to L$ are contiguous, meaning that the union of their vertex
images of each source simplex is a target simplex, then
\[
 P[v_0,\ldots,v_k]
 =\sum_{i=0}^{k}(-1)^i
 [f(v_0),\ldots,f(v_i),g(v_i),\ldots,g(v_k)]
\]
is a finite chain homotopy satisfying
$\partial P+P\partial=g_\#-f_\#$.
\end{proposition}
\begin{proof}
The alternating face sum is well defined under the oriented-tuple
conventions.  For a tuple with a repeated vertex, the two terms deleting
one of those repetitions cancel after orienting the remaining tuples;
all other terms still have a repetition and vanish.  The face sum also
respects the sign under an adjacent exchange, as checked by pairing the
two exchanged-index terms and reindexing the others.  Adjacent exchanges
generate every reordering.  The chain-map identity therefore follows
by applying $f$ to each term of the face sum, whether its image is
nondegenerate or degenerate.

To verify the prism identity, expand the two finite double sums.
Deleting an $f(v_j)$ with $j<i$, or a $g(v_j)$ with $j>i$, pairs with the
prism summand for the face omitting $v_j$; their signs are opposite on
the two sides of $\partial P+P\partial$.  Among the remaining terms,
deleting the last $f$-vertex in one summand cancels deleting the first
$g$-vertex in the preceding summand.  The first and last unpaired terms
are respectively $[g(v_0),\ldots,g(v_k)]$ and
$-[f(v_0),\ldots,f(v_k)]$.  This proves the identity.  The same
pairings remain valid when some terms vanish by repetition.  Contiguity
ensures that every displayed tuple belongs to $L$.  All index ranges
are bounded by $k+1$, so the calculation is a finite integer identity.
\end{proof}

\begin{theorem}[A20: an explicit simplex-carrier homotopy]
\label{thm:simplex-carrier}
Assign to each simplex $\sigma$ of a finite complex $K$ a nonempty
simplex $C(\sigma)$ of $L$, with
$C(\rho)\subseteq C(\sigma)$ whenever $\rho$ is a nonempty face of
$\sigma$.  Suppose two supplied chain maps $F,G$ have their value on
$\sigma$ supported in $C(\sigma)$ and preserve augmentation in degree
zero.  There is a primitive recursively constructed chain homotopy
between them, also supported in these carriers.
\end{theorem}
\begin{proof}
Augment each simplex chain complex by $C_{-1}=\Q$ and
$\partial[v]=1$.  For a chosen vertex $a$ of a simplex, let
$c_a[w_0,\ldots,w_k]=[a,w_0,\ldots,w_k]$ and $c_a(1)=[a]$.
Deleting the first vertex gives the original tuple; every other face
cancels its counterpart in $c_a\partial$.  Hence
$\partial c_a+c_a\partial=1$, including degree zero.  Choose for
$a_\sigma$ the least vertex of $C(\sigma)$.

Define $H_{-1}=0$.  After $H$ has been defined on proper faces of
$\sigma$, put
\[
 z_\sigma=G(\sigma)-F(\sigma)-H(\partial\sigma),\qquad
 H(\sigma)=c_{a_\sigma}z_\sigma.
\]
The induction hypothesis on the faces, the chain-map identities, and
$\partial^2=0$ give
\[
 \partial z_\sigma
   =(G-F)(\partial\sigma)-\partial H(\partial\sigma)=0.
\]
For a vertex this says that $G(v)-F(v)$ has augmentation zero, precisely
the stated hypothesis.  Nested carriers ensure that all terms of
$z_\sigma$ lie in $C(\sigma)$.  The cone identity gives
$\partial H(\sigma)=z_\sigma$, which rearranges to
$G-F=\partial H+H\partial$.  This completes induction in increasing
simplex dimension.  Only faces already processed are consulted; the
number of simplices and their dimensions bound this recursion.  Thus
the construction consists entirely of bounded list operations and
rational additions and yields the promised explicit certificate.
\end{proof}

\subsection{Graphs and surface presentations}

\begin{theorem}[A07: graph Betti numbers and a cycle basis]
\label{thm:graph-betti}
Let a finite loopless undirected graph have $v$ vertices, $e$ edges, and
$c$ connected components, with $c=0$ for the empty graph.  Viewed as a
one-dimensional chain complex over $\Q$,
\[
 b_0=c,\qquad b_1=e-v+c,\qquad \chi=v-e.
\]
A spanning forest and a basis of $b_1$ cycles are primitive recursively
computable from the adjacency table.
\end{theorem}
\begin{proof}
Start from the least unvisited vertex; repeatedly add the least edge with
exactly one endpoint visited until none remains, then start the next
component.  At most $v$ vertices are added and every edge scan has length
$e$.  The selected edges form a tree in each component.  Orient all
edges, and assign to $u\to w$ the boundary $w-u$.  In a component with
root $r$, the boundaries of the tree paths from $r$ to its other vertices
span the vectors $w-r$.  These $v_i-1$ vectors are independent by their
nonroot coordinates and span exactly the vectors whose coordinate sum on
that component is zero.  Every edge boundary has coordinate sum zero.
Thus the boundary rank is $v-c$.

For each nontree edge, take that oriented edge minus the oriented tree
path with the same endpoints.  Its boundary is zero.  These cycles are
independent because each has coefficient $1$ on its own nontree edge
and coefficient $0$ on the others.  Subtracting their suitable multiples
from any cycle leaves a cycle supported on the forest.  A leaf of a
nontrivial tree forces the coefficient of its unique incident edge to
vanish; removing leaves by at most $v$ bounded stages shows that a forest
has no nonzero cycles.  Hence these $e-(v-c)$ cycles form a basis.
The stated formulas follow, including the zero-size case.
\end{proof}

\begin{proposition}[A08: genus formulas for supplied polygonal models]
\label{prop:polygon-genus}
For the orientable polygonal cellular presentation with one vertex,
$2g$ oriented edges, and one two-cell attached along
$\prod_{i=1}^{g}a_i b_i a_i^{-1}b_i^{-1}$, its rational cellular chain
complex has
\[
 (b_0,b_1,b_2)=(1,2g,1),\qquad \chi=2-2g.
\]
For the nonorientable presentation with one vertex, $h\geq1$ edges, and
attaching word $a_1^2\cdots a_h^2$,
\[
 (b_0,b_1,b_2)=(1,h-1,0),\qquad \chi=2-h.
\]
\end{proposition}
\begin{proof}
Every edge begins and ends at the same vertex, so $D_1=0$.  In the
specified cellular complex, the coefficient of an edge in $D_2$ is its
signed number of traversals in the attaching word.  The commutator word
has coefficient zero on every edge, giving $D_2=0$.  The square word gives
the nonzero column $(2,\ldots,2)^t$, of rank one over $\Q$.
Proposition~\ref{prop:effective-homology} gives the formulas.  At $g=0$
the presentation has no edges and a two-cell attached to the sole vertex,
giving the indicated sphere chain complex.
\end{proof}

The theorem calculates the \emph{supplied} cellular presentation.  To
transfer it to a particular finite complex one may supply a chain-homotopy
equivalence and apply Theorem~\ref{thm:chain-homotopy}.  No assertion is
made here that an arbitrary surface can be recognized, classified, or
reduced to one of these presentations.  In particular a table of homology
matrices alone does not certify that its realization is a surface.

\begin{proposition}[A09: a bounded coloring theorem]
A loopless graph of maximum degree $\Delta$ has a proper coloring with
$\Delta+1$ colors.  A coloring, and the least number of colors admitting a
coloring, can be found by primitive recursive bounded search.
\end{proposition}
\begin{proof}
Process vertices in order.  The already colored neighbors use at most
$\Delta$ colors, so scanning $0,\ldots,\Delta$ finds an unused color.
At the step when the later endpoint of any edge is processed, it receives
a different color from the earlier endpoint.  This proves correctness of
the greedy construction by induction on the vertex index.  To find the
chromatic number, test all $q^v$ assignments for $q=0,1,\ldots,v$; the
unique assignment for the empty graph is handled separately.  The
$v$-color assignment exists by giving distinct colors to distinct
vertices, so the bounded search succeeds.  This is not a proof of the
four-color theorem mentioned in the source entry.
\end{proof}

\begin{theorem}[A10: two colors or an odd-cycle certificate]
A finite loopless graph admits a two-coloring if and only if it has no
odd cycle.  A bounded algorithm returns either a two-coloring or a simple
odd cycle of length at most $v$.
\end{theorem}
\begin{proof}
Construct a spanning forest by breadth-first search and color each vertex
by the parity of its tree distance from the root.  If every edge joins
opposite parities, this is the required coloring.  Otherwise take an edge
$u w$ with equal endpoint parities.  The unique tree path from $u$ to
$w$ has even length: removing the common initial segment of their root
paths removes twice the same number of edges.  The offending edge cannot
be a tree edge.  Appending it to that simple tree path gives a simple odd
cycle, using at most $v$ vertices.  Conversely, colors alternate around
every cycle in a proper two-coloring; returning to the original color
requires an even number of edges.  Every search and path extraction is
bounded by the finite graph size.
\end{proof}

\subsection{Finite sheaves with executable gluing}

Fix a finite topological space $X=\{0,\ldots,n-1\}$ with its open subsets
listed as bit masks.  Let $N_x$ be the intersection of its open
neighborhoods.  This is open, contains $x$, and is contained in every
open set containing $x$.  A finite-set-valued presheaf is a table of
finite sets $F(U)=\{0,\ldots,m_U-1\}$ and restriction tables
\[
 \rho_{U,V}:F(U)\longrightarrow F(V)\quad(V\subseteq U)
\]
satisfying $\rho_{U,U}=1$ and
$\rho_{V,W}\rho_{U,V}=\rho_{U,W}$.  These identities are bounded tests.
A matching family on a listed finite cover $(U_i)_{i<\ell}$ of $U$
consists of $s_i\in F(U_i)$ with equal restrictions to every
$U_i\cap U_j$.  A gluing is $s\in F(U)$ restricting to all $s_i$.

\begin{theorem}[A11: canonical-cover test for a finite sheaf]
\label{thm:canonical-sheaf}
Such a presheaf has unique gluing on every finite open cover if and only
if it has unique gluing on the covers $(N_x)_{x\in U}$ for every open
$U$.  The $U=\varnothing$ case requires $F(\varnothing)$ to be a
singleton.  Both this property and the resulting gluing operation are
primitive recursively decidable from the tables.
\end{theorem}
\begin{proof}
One implication is immediate by using these particular covers.  For the
other, let $(U_i,s_i)$ be an arbitrary finite matching family.  For each
$x\in U$, select the least index $i(x)$ with $x\in U_{i(x)}$, and put
$t_x=s_{i(x)}|_{N_x}$.  If $x,y\in U$, their two original sections agree
on $U_{i(x)}\cap U_{i(y)}$, hence $t_x,t_y$ agree on
$N_x\cap N_y$ by restriction transitivity.  Canonical gluing supplies
$s\in F(U)$.  For every $i$ and $x\in U_i$, compatibility of the original
family gives $s|_{N_x}=s_i|_{N_x}$.  Canonical uniqueness on $U_i$
therefore gives $s|_{U_i}=s_i$.  Two possible global gluings agree on
every $N_x$, so canonical uniqueness on $U$ gives equality.

For the algorithm, let $M=1+\max_U m_U$.  There are at most $2^n$ open
sets, at most $M^n$ tuples on any canonical cover, and at most $M$
candidate global sections.  Test compatibility, then count global
sections whose restrictions have the prescribed values; require this
count to be exactly one.  This bounds all tests.  Once the test succeeds,
searching the $m_U$ entries of $F(U)$ returns the unique gluing.
\end{proof}

\begin{proposition}[A12: function sheaves, continuity, and rings]
For a finite set $A$, the assignment $U\mapsto A^U$ is a finite sheaf.
For a finite topological space $Y$, its subassignment consisting of
continuous maps $U\to Y$ is also a sheaf.  If $R$ is a finite ring with
discrete topology, locally constant $R$-valued maps form a sheaf of rings.
\end{proposition}
\begin{proof}
For matching local functions $f_i$, define $f(x)=f_{i(x)}(x)$ using the
least covering index.  Agreement on overlaps proves that this value is
independent of the covering member, and proves that $f|_{U_i}=f_i$.
Any other gluing has the same value at each point, proving uniqueness.
For an open subset $V$ of $Y$,
$f^{-1}(V)=\bigcup_i f_i^{-1}(V)$.  Each inverse image on the right is
open in $U_i$, hence in $U$, and their finite union is open.  Thus the
gluing is continuous.  Into a discrete finite space, continuity is
equivalent to local constancy: fibers are open, and conversely the open
fibers give inverse images of all subsets.  Pointwise addition, negation,
and multiplication of locally constant maps remain locally constant
by intersecting their finitely many relevant neighborhoods.  Ring laws
hold at each point, and restrictions preserve these operations and the
unit.  This gives the ring-valued assertion, including the one-element
section ring on the empty set.
\end{proof}

\begin{proposition}[A13: restriction preserves the sheaf property]
If $W$ is open in $X$, define $(F|_W)(U)=F(U)$ for open $U\subseteq W$,
using the same restrictions.  A finite sheaf restricts to a finite sheaf,
and repeated restrictions agree exactly.
\end{proposition}
\begin{proof}
Since $W$ is open, every open subset of $W$ is open in $X$.  A cover,
matching family, or gluing in $W$ is consequently the same finite data
in $X$.  Existence and uniqueness are inherited unchanged.  Restricting
again merely deletes more rows from the same tables, which proves the
composition assertion.
\end{proof}

\begin{proposition}[A14: finite stalks and local equality]
\label{prop:finite-stalk}
For any finite presheaf, its stalk at $x$ has the canonical model
$F_x=F(N_x)$.  Two local sections have the same germ at $x$ precisely when
their restrictions to $N_x$ agree.  For a sheaf and $s,t\in F(U)$,
\[
 s=t\quad\Longleftrightarrow\quad
       (\forall x\in U)\ s|_{N_x}=t|_{N_x}.
\]
\end{proposition}
\begin{proof}
Every neighborhood $V$ of $x$ contains $N_x$, giving the map
$(V,s)\mapsto s|_{N_x}$.  If two representatives agree on some smaller
neighborhood of $x$, they agree after restricting further to $N_x$.
Conversely, equality on $N_x$ itself witnesses equality of germs.
Every element of $F(N_x)$ is represented by $(N_x,s)$, proving the
claimed model and its inverse.  The last assertion is exactly uniqueness
on the canonical cover of $U$.
\end{proof}

\begin{theorem}[A15: stalkwise isomorphisms of finite sheaves]
Let $\alpha:F\to G$ be tables commuting with restrictions between finite
sheaves.  If every $\alpha_{N_x}$ is bijective, then every $\alpha_U$ is
bijective.  An inverse natural transformation is primitive recursively
computable.
\end{theorem}
\begin{proof}
If $\alpha_U(s)=\alpha_U(t)$, restrict to each $N_x$ and use injectivity
there.  Proposition~\ref{prop:finite-stalk} gives $s=t$; hence every
$\alpha_U$ is injective.  Given $t\in G(U)$, use bounded search for the
unique $s_x\in F(N_x)$ mapping to $t|_{N_x}$.  On
$V=N_x\cap N_y$, the images of $s_x|_V$ and $s_y|_V$ agree.  The
already proved injectivity of $\alpha_V$ gives agreement of the two
restrictions.  Glue the $s_x$ to $s\in F(U)$.  Its image equals $t$ on
all $N_x$, and hence equals $t$.  This proves surjectivity, including
$U=\varnothing$ because both sheaves have one empty section.
The inverse commutes with restrictions by applying $\alpha$ and using
its injectivity.  All searches have lengths bounded by the section
tables.
\end{proof}

\begin{proposition}[A16: local subobjects and equalizer sheaves]
Suppose $F$ is a finite sheaf and $G(U)\subseteq F(U)$ are decidable
subsets preserved by restrictions.  Assume locality: if a section
$s\in F(U)$ has its restrictions in $G(U_i)$ on a finite cover, then
$s\in G(U)$.  Then $G$ is a sheaf.  In particular, for two morphisms
$\alpha,\beta:F\to H$ between finite sheaves, the tables
$G(U)=\{s:\alpha_U(s)=\beta_U(s)\}$ form a subsheaf.  For ring
morphisms these are subrings.
\end{proposition}
\begin{proof}
A matching family in $G$ has a unique gluing in $F$; locality puts that
gluing in $G$.  Uniqueness is inherited from $F$.  In the equalizer
case, commuting with restrictions gives restriction stability.  If
$\alpha(s)$ and $\beta(s)$ agree on a cover, uniqueness in $H$ gives
their global equality, proving locality.  Preservation of ring operations
by $\alpha$ and $\beta$ proves closure under those operations.
\end{proof}

\paragraph{Presheaves of categories: precise limited coverage.}
Finite categories can be encoded by object and arrow lists, source,
target, identity, and composition tables.  A finite diagram of restriction
functors is checked by bounded tests of these tables.  For the
category-valued entry's more general pseudofunctor convention, include
the components of its natural isomorphisms, their inverse tables, and
the coherence equations; each displayed equation is again a bounded
test.  This supplies definitions and a checker interface for a
\emph{specified finite diagram}.  We do not infer descent for categories
from the set-valued sheaf theorem, nor construct all local homeomorphisms
of an unrestricted space.  Those would be additional results.

\subsection{Finite Boolean algebras and the countable-union boundary}

\begin{theorem}[A17: generated Boolean algebras by membership atoms]
\label{thm:boolean-atoms}
Given subsets $E_0,\ldots,E_{s-1}$ of a finite $n$-element set, define
$x\equiv y$ when their $s$ membership bits agree.  The nonempty
equivalence classes are the atoms of the generated Boolean algebra.
The generated algebra is exactly the collection of unions of these
atoms; if there are $a$ atoms, it has $2^a$ members, with
$a\leq\min(n,2^s)$.  Its complete table is primitive recursively
computable.
\end{theorem}
\begin{proof}
Compute the membership word of every point.  For each word $w<2^s$,
the corresponding class is
\[
 A_w=\bigcap_{i<s}
 \begin{cases}E_i,&w_i=1,\\ X\setminus E_i,&w_i=0.\end{cases}
\]
The empty intersection is $X$.  Thus every nonempty class belongs to
every Boolean algebra containing the generators.  Unions of the
nonempty classes are closed under union and complement, because the
classes partition $X$, and contain every generator.  They are therefore
exactly the smallest Boolean algebra containing the generators.
Different subsets of the nonempty classes have different unions: choose
the least point in a class appearing in just one of the subsets.  This
gives the cardinality $2^a$.  Each class contains a different least point
and has a different word, giving both bounds for $a$.  Enumerating the
$2^s$ words and $2^a$ unions supplies explicit bounded algorithms.
\end{proof}

\begin{proposition}[A18: bounded and stabilization-certified unions]
For such an algebra, the union of any listed finite family is a member
and its mask is computed by bitwise disjunction.  For a fixed PR sequence
$E(t)$ of algebra members, if a number $T$ and the property
\[
   \forall t\quad E(t)\subseteq\bigcup_{i<T}E(i)
\]
are supplied, the finite union on the right represents the unbounded
union predicate $x\mapsto\exists t\,(x\in E(t))$.
\end{proposition}
\begin{proof}
Finite union closure follows either from its Boolean operations or from
the atom representation.  Under the supplied property, a membership
witness $t$ for the unbounded union implies membership in the finite
union.  Conversely, membership in a finite union gives an index $i<T$
by bounded search, hence a witness for the existential predicate.  These
are intuitionistically valid implications.  The displayed universal
property is a hypothesis; a numerical cutoff alone does not certify it.
\end{proof}

\paragraph{Why this is not automatic countable-union closure.}
An arbitrary unbounded union of coded finite sets need not have a
\emph{uniformly computable Boolean code}.  On $X=\{0\}$ let $E_e(t)$
be $\{0\}$ when machine $e$ on input $e$ has halted within $t$ steps and empty
otherwise.  This is uniformly primitive recursive.  A uniform algorithm
returning the one-bit code of $\bigcup_t E_e(t)$ would decide halting by
reading its sole bit.  Indeed, if that total decider were $D$, the program
which on input $j$ halts exactly when $D(j)=0$ has some index $p$.
Running it on $p$ contradicts either answer of $D(p)$.  Thus finiteness of the underlying set does not
provide a uniform terminating union algorithm.  A monotone chain can
change at most $n$ times while its last change still has no computable
bound.  The theorem above is the exact replacement used here for the
\emph{Sigma-Algebra} entry: generated finite Boolean algebras, finite
unions, and unions with an explicitly justified stabilization bound.

\subsection{Formalization and the arithmetic proof boundary}

The file \texttt{AlgebraicBlueprint.lean} gives executable rational
matrix operations and precise propositions for finite chain-complex
certificates.  It includes checked structural lemmas for integral-chain
homotopies on cycles, for gluing supplied compatible functions, and for
the full Euler telescoping identity from supplied rank--nullity data.
Uncompleted matrix-rank and finite-sheaf
theorems are named propositions with their required data, not axioms or
theorems with hidden proof holes.

For a complete formalization, first prove the rational-pair arithmetic
and bounded row-operation invariants, then the cycle/boundary basis
algorithm of Proposition~\ref{prop:effective-homology}.  The
Euler--Poincar\'e theorem is then a finite telescoping identity; chain
homotopies and collapses are matrix identities.  Independently, encode
open and section tables by bounded lists and implement the canonical-cover
checker of Theorem~\ref{thm:canonical-sheaf}.  Its proof only uses bounded
selection and the table equations.  The graph and atom constructions
use bounded graph searches and bit-mask loops.

For the claimed \ipra\ interpretation, the sizes, loops, rational-height
bounds, and matrix identities must be re-expressed with the chosen
primitive recursive coding and quantifier-free induction.  A universal
quantifier over a finite list is a Boolean fold with an explicit bound;
its induction invariant is a Boolean equality about the current prefix.
The human proofs above specify those loops and invariants, and therefore
give a direct route.  Checked Lean proofs about functions, matrices, or
finite types in the ambient metatheory are not themselves an
object-language derivation in \ipra.  The supplied tests likewise check
examples and implementations, not universal arithmetic provability.
